\section{Introduction}
\label{sec:intro}

Venture-backed governance operates through a three-tier architecture that has received remarkably asymmetric academic attention. Directors---with fiduciary duties, voting rights, and personal liability---are the subject of a vast literature \citep{hermalin2003boards, adams2010role}. Board advisors occupy an informal tier with limited access. Between them sit board observers: individuals who attend every board meeting, receive every board document, and participate in every strategic discussion, yet owe no fiduciary duty to the company, hold no voting rights, and bear no personal liability. According to the NVCA's Q4 2023 CFO Working-Group Survey, 82\% of VC entities appoint board observers, with universal adoption among firms managing more than \$500 million in assets \citep{packin2025board}. Despite this prevalence, not a single empirical business paper has studied them. \citet{carter2025} calls explicitly for research on board observers, and \citet{ewens2024board} acknowledge their existence in a footnote before excluding them from analysis.

This paper asks whether the observer role creates information channels that would not exist under a traditional director-only board. Observers possess the same information as directors---they hear the same discussions about acquisitions, financing, restructuring, and personnel decisions---but face none of the legal constraints that govern how directors handle that information. We term this asymmetry the \textit{accountability gap}. If this gap matters for information flow, the effects should be detectable in the stock prices of public companies connected to observers through their concurrent board roles.

We construct a person-level network linking observers at private firms to their directorships at public firms. The same individual who sits as a nonvoting observer at a private company also serves as a voting director at a public company, personally bridging the two boardrooms. Using 57,000 material events at private companies from S\&P Capital IQ Key Developments (2015--2025) and daily stock returns from CRSP, we test whether connected public firms show abnormal returns in the weeks before events at the private firms become public.

Three sets of findings emerge. First, connected portfolio companies earn higher market-adjusted returns than non-connected stocks around the same events. For executive and board changes---the most common event type---connected firms outperform by 0.28 percentage points in the five days before the announcement ($p=0.010$, event-clustered standard errors), a result that survives six alternative specifications.

Second, the effect is concentrated in events where the board discusses something material and confidential, and where the information is industry-relevant. When an observer's private company announces an acquisition, same-industry connected firms earn cumulative abnormal returns of 4.6 percentage points over the prior month ($p=0.001$). Around bankruptcy filings, the figure is 9.5 percentage points ($p=0.043$). When the observed company is being acquired---the tightest of board secrets---the effect appears only at the moment of announcement: 1.1 percentage points at day $[-1,0]$ ($p=0.027$). For routine events such as private placements and product announcements, we find no effect. Neither the connection alone nor the industry match alone produces the result; only the interaction of both---being connected through the observer network \textit{and} operating in the same industry---generates abnormal returns.

Third, the information channel responds to regulatory changes affecting observer accountability. In 2020, the NVCA removed fiduciary language from its standard observer provisions, loosening the perceived obligation of observers to safeguard private information. Same-industry spillover, previously indistinguishable from zero, became significant ($p=0.001$). In January 2025, the DOJ and FTC explicitly extended Clayton Act Section~8 interlocking directorate rules to board observers for the first time, subjecting same-industry connections to antitrust scrutiny. Same-industry spillover reversed ($p<0.001$). Two independent regulatory changes, operating in opposite directions on the same mechanism, with clean placebo tests at alternative break dates.

Our contribution is threefold. First, we document a private-to-public information channel that prior work on interlocking directorates has not examined. The existing literature studies spillovers between public firms connected through shared directors \citep{kang2008director, shi2025shared, fracassi2017corporate}, where both firms' information is semi-public and Regulation Fair Disclosure constrains selective disclosure. Our setting differs: the observed firm is private, its board meetings are exempt from Reg~FD, and the observer---who owes no fiduciary duty to the private firm---carries the information to a public firm where it reaches prices. Second, we identify \textit{pre-event} information leakage rather than post-event reputation spillovers or decision mimicry. Connected firms move before the event becomes public, not after. Third, we exploit two quasi-experimental regulatory shocks---the NVCA template change and the Clayton Act extension---that directly alter the information constraints on observers, providing causal variation that complements our cross-sectional evidence.

We frame our findings as evidence of \textit{information permeability} of the observer network rather than insider trading. We show that information flows through the network and reaches prices, but we do not claim to identify the specific transmission mechanism. The stock price effects could arise from trading by the observer or their associates, from strategic decisions at the connected public firm informed by the observer's dual role, or from broader information diffusion through the VC ecosystem. The regulatory implications are direct regardless of mechanism: the DOJ/FTC's extension of antitrust scrutiny to board observers is empirically supported by our data.

The remainder of the paper is organized as follows. Section~\ref{sec:institutional} describes the institutional setting of board observers, including the relevant regulatory framework. Section~\ref{sec:data} presents the data and research design. Section~\ref{sec:results} reports our findings. Section~\ref{sec:conclusion} concludes.